\section{The graphic, sections, and complete displacement}\label{sec:geometry}
All estimates below are uniform on a sufficiently small compact interval
$I_b\Subset(0,2)$ about a prescribed $b_*$.

We use the physical displacement \eqref{eq:common-displacement} from
Section~\ref{sec:return-framework}. The table records the regional
coordinates and domains before they are transferred to those sections.

\begin{table}[H]
\centering
\caption{Sections, coordinates, maps, and actual domains.}
\label{tab:return-objects}
\small
\setlength{\tabcolsep}{3pt}
\begin{tabular}{@{}>{\raggedright\arraybackslash}p{0.14\linewidth}
                >{\raggedright\arraybackslash}p{0.27\linewidth}
                >{\raggedright\arraybackslash}p{0.18\linewidth}
                >{\raggedright\arraybackslash}p{0.32\linewidth}@{}}
\toprule
Object & Source and target & Map & Actual domain and orientation\\
\midrule
Regular branch & $\Sigma_u(u_{\rm up})$ to $\Sigma_\xi(E_b)$
 & $T_-$ & fixed regular core; $(T_-)_{u_{\rm up}}=-j<0$\\
Contact branch & $\Sigma_u(u_{\rm up})$ to $\Sigma_\xi(E_b)$
 & $T_+$ & open finite-passage components common with $T_-$\\
Displacement & same source and target
 & $\Delta=T_+-T_-$ & zeros are periodic orbits; isolated zeros are limit cycles\\
Negative passage & prescribed physical source to $z=-Ln$
 & $J_-$ & positive right-going finite passage; local source coordinate $u$ is not $u_{\rm up}$\\
Positive complement & positive-saddle outgoing cut to $z=z_c$
 & $C_+$ & relative neighborhood of a same-field positive reference; oriented so $C_+'>0$\\
Saddle collars & fixed incoming to fixed outgoing cuts
 & $\mathcal D$ & open finite-time side; a zero-amplitude separatrix is a boundary\\
Root collision & left base to through-passage or fiber-exit target
 & branch maps & localized actual interval $0<A<\eta_0$; at most one finite event value\\
Lower connector & lower power output to the positive-saddle incoming cut
 & $\Psi$ & increasing physical passage interval;
 $\Psi(0)=-B_{\rm conn}<0$ is retained\\
\bottomrule
\end{tabular}
\end{table}

\subsection{Fast orbits and transverse sections}
We first identify the limiting fast arc and then choose sections that
persist for the perturbed fields. This supplies the geometric input for
all subsequent transition maps.
Set $\omega_b=\sqrt{4-b^2}/2$. The solution of~\eqref{eq:fast}
through the origin is
\begin{equation}\label{eq:fast-orbit}
 z(t)=\frac{e^{bt/2}}{\omega_b}\sin(\omega_bt),\qquad
 y(t)=1+e^{bt/2}\left(-\cos(\omega_bt)
                  +\frac{b}{2\omega_b}\sin(\omega_bt)\right).
\end{equation}
The ordinate increases until $t=\pi/\omega_b$ and then decreases
until $2\pi/\omega_b$, where it is negative. Its first positive
zero $t_b$ is therefore unique and belongs to
$(\pi/\omega_b,2\pi/\omega_b)$. At that zero $z(t_b)<0$.
The implicit function theorem shows that $f_b=z(t_b)$ and the whole
arc depend smoothly on $b$. In particular $-f_b$ has a positive
lower bound on $I_b$. The arc is simple: it is part of the trajectory
of the linear focus~\eqref{eq:fast}, before its first return to the
horizontal axis. Thus~\eqref{eq:graphic} is a Jordan curve.
Continuous dependence, together with the transverse return at $f_b$,
also proves $\Gamma_x^b\to\Gamma_0^b$ in Hausdorff distance as
$x\downarrow0$, uniformly on $I_b$.

The upper intersection of the arc with $z=0$ has height
\[
 y_b=1+\exp\left(\frac{\pi b}{\sqrt{4-b^2}}\right).
\]
We use the section
$\Sigma_u=\{(0,y_b+u_{\rm up}):|u_{\rm up}|<u_0\}$ and a
small outgoing section $\Sigma_\xi$ on $z=\xi>0$, near the lower
intersection with the fast arc. The coordinate $u_{\rm up}$ ranges over a
fixed interval containing zero, independently of every later rescaling.
The endpoint parameter $x$ indexes the unperturbed family of graphics.
At the contact endpoint $x=0$, every regional construction is pulled back
to the regular physical coordinate $u_{\rm up}$; $x$ is not a return
coordinate.

\begin{lemma}\label{lem:fast-energy}
On a fixed rectangle about the contact point where $1+bz-y>0$,
there is an analytic first integral $E_b$ of~\eqref{eq:fast} such that
\begin{equation}\label{eq:energy-fast}
 (E_b)_y>0,\qquad
 (E_b)_z+\frac{z}{1+bz-y}(E_b)_y=0,\qquad
 E_b(0,y)=y.
\end{equation}
Its $y$ derivative is a positive unit on a smaller rectangle. The
transition along any fixed compact part of the fast arc, written in
this first-integral coordinate at its two ends, converges to the
identity in $C^k$ on the prescribed compact source interval, uniformly
for $b\in I_b$, for every fixed finite $k$, as $\mu\to0$.
\end{lemma}
\begin{proof}
Since the $z$ component of~\eqref{eq:fast} is nonzero there, each
orbit meets $z=0$ once locally. Define $E_b$ to be the ordinate of
that intersection. The analytic dependence theorem gives
\eqref{eq:energy-fast} and $(E_b)_y(0,0)=1$. Continue this first
integral along the required regular fast arc. The perturbed transition
between fixed transverse sections depends smoothly on the parameters;
at $\mu=0$ its energy coordinate is constant. This proves the last
assertion. Sections must be fixed before taking this limit.
\end{proof}

For $\mu\ne0$ we shall write the family exactly as
\begin{equation}\label{eq:normalized}
 F:=\dot z=y(1+bz-y)+\varepsilon^2P(z),\qquad
 G:=\dot y=zy+\varepsilon^3D,
 \qquad P(z)=E_0+E_1z+E_2z^2,
\end{equation}
with bounded coefficients. The normalization covering every raw
parameter is defined in Section~\ref{sec:return-framework};
Section~\ref{sec:parameter-cover} completes its finite directional cover.

\begin{lemma}\label{lem:sections}
After shrinking the parameter and Hausdorff neighborhoods, every
cycle counted in Theorem~\ref{thm:main} lies in $y>0$ and meets
$\Sigma_u$ and every member of any prescribed finite collection of
regular fast sections exactly once, inside prescribed smaller section
intervals. The backward transition $T_-$ from $\Sigma_u$ to
$\Sigma_\xi$, measured by $E_b$, satisfies
\begin{equation}\label{eq:short-map}
 (T_-)_{u_{\rm up}}=-j(u_{\rm up},p),\qquad c\le j\le C,
 \qquad |\partial_{u_{\rm up}}\log j|\le C.
\end{equation}
Here and below $p$ denotes the fixed parameters of a transition.

Moreover, if $P_p$ is the transition between two fixed vertical sections
contained in a compact regular flow tube on the fast arc, then, after one
common shrinking of the parameter neighborhood,
\begin{equation}\label{eq:regular-holonomy-jets}
 0<c\le |P_p'|\le C,\qquad
 \left|\frac{P_p''}{P_p'}\right|
 +\left|\left(\frac{P_p''}{P_p'}\right)'\right|\le C .
\end{equation}
If the transverse speeds on the two sections have the same fixed sign,
as for the holonomies on the outgoing regular arc, then $P_p'>0$.
The same assertion holds when either section is measured in the energy
coordinate $E_b$.

For bounded coefficients in~\eqref{eq:normalized}, the forward orbit
from the same signed section reaches $y=K\varepsilon^3$ at
$z=\Lambda(u_{\rm up},p)$ in a fixed negative interval about $f_b$, where
\begin{equation}\label{eq:relaxation-map}
 -C\le\Lambda_{u_{\rm up}}\le-c<0,\qquad
 |\Lambda_{u_{\rm up}u_{\rm up}}|\le C.
\end{equation}
Suppose in addition that $P\ge p_0>0$ on a fixed interval
$[z_{\rm ent},-q_g]$, where $z_{\rm ent}<-q_g$, from these entries to
$-q_g<0$. Then every counted cycle reaches
$z=-q_g$ at $0<y\le K_1\varepsilon^3$, and its subsequent arc
to $\Sigma_\xi$ has $F>0$. The union of admissible passages is
represented on one open interval of $u_{\rm up}$ by a forward transition $T_+$
to the same target. Its intersections with $T_-$ are precisely the
periodic orbits obtained by finite passages there; isolated intersections
represent limit cycles.
\end{lemma}
\begin{proof}
At each regular marker the organizing field is nonzero and transverse
to its section. A uniformly small flow box implies that any orbit
meeting a smaller ball also meets the section inside its prescribed
core. One Hausdorff inclusion provides such a point on every counted
cycle. Choose the sections so that their portions outside the prescribed
cores have positive distance from the graphic. The other inclusion then
confines every intersection to those cores.
A periodic orbit is a Jordan curve. As a transverse arc is traversed,
its intersections with that curve alternate between entering and
leaving the bounded component. Their oriented intersection signs
alternate, whereas the field has a fixed transverse sign on the
section. Hence there can be only one intersection.

One marker has a fixed positive ordinate. If $\mu_4\ne0$, every
crossing of $y=0$ has the same sign, which is impossible for a
periodic ordinate. If $\mu_4=0$, that line is invariant; a cycle close
to the graphic cannot cross it, and Hausdorff closeness to the positive
fast arc excludes a cycle in $y<0$. Thus the entire cycle lies in
$y>0$. The compact regular arc from
$\Sigma_\xi$ to $\Sigma_u$ persists. Fix a closed flow tube around this
arc, with slightly larger section cores. On the limiting field the tube
contains no zero and the transverse speeds on its vertical sections are
bounded away from zero. Compactness in $b$ and continuity give the same
bounds, a common flight-time bound, and common $C^4$ bounds for the
normalized field after shrinking the parameter neighborhood. The target
event equation has time derivative bounded below, so the implicit-function
theorem and the first three variational equations give a common $C^3$ bound
for every such holonomy map.

Writing $X_p=(F_p,G_p)$ for the normalized vector field, a transition
$P_p(y_0)=y_1$ between vertical sections satisfies the planar determinant
identity
\begin{equation}\label{eq:holonomy-determinant}
 F_p(z_1,y_1)P_p'(y_0)=F_p(z_0,y_0)
 \exp\!\left(\int_0^{T(y_0,p)}\operatorname{div}X_p(\phi_t)\,dt\right).
\end{equation}
The event-time correction is parallel to the target vector field and does
not alter this identity. The uniform speed and flight-time bounds give a
positive lower bound for $|P_p'|$ as well as an upper bound. On the outgoing
regular arc the two transverse speeds have the same fixed sign, hence
$P_p'>0$. Combining these facts with the $C^3$ estimate yields
~\eqref{eq:regular-holonomy-jets}.
In the energy coordinate, the section change has derivative $(E_b)_y$, a
positive $C^3$ unit on the fixed tube, so the composition rules preserve
the same bounds. Taking the two sections to be $\Sigma_u$ and $\Sigma_\xi$,
the signs $F<0$ on $\Sigma_u$ and $F>0$ on $\Sigma_\xi$ give the negative
orientation $(T_-)_{u_{\rm up}}=-j(u_{\rm up},p)$ with $j$ a positive
unit, as in
~\eqref{eq:short-map}. Its limiting energy derivative does not depend on
the final position of $\Sigma_\xi$ on that arc.

Follow the descending regular arc to a fixed small height $y_h$.
For $K\varepsilon^3\le y\le y_h$ and $z$ in the negative entry
interval, choose $K$ so that $zy+\varepsilon^3D\le-cy$.
With $y$ as independent variable, the difference between the perturbed
and divided fast equations is exactly
\begin{equation}\label{eq:relaxation-error}
 \frac{\varepsilon^2zP(z)
       -\varepsilon^3D(1+bz-y)}{z(zy+\varepsilon^3D)}.
\end{equation}
This function and its first two $z$ derivatives are bounded by
$C\varepsilon^2/y$. Integration and the first two variational
equations give a $C^2$ error of size
$O(\varepsilon^2(1+|\log\varepsilon|))$ in the entry abscissa.
The unperturbed map from the upper section to the lower fast foot
has negative derivative bounded away from zero, by the determinant
formula on the divided regular field. This proves
\eqref{eq:relaxation-map} without a sign assumption on $P$.

Under the additional positive-polynomial hypothesis, put
$h_s=-\varepsilon^3D/z$ and $w=y-h_s$. Then
\begin{align*}
 F&=\varepsilon^2P+h_s(1+bz-h_s)
                  +(1+bz-2h_s)w-w^2,\\
 \dot w&=zw-h_s'F.
\end{align*}
On $|w|\le K_1\varepsilon^3$, the horizontal speed is
$\frac12p_0\varepsilon^2\le F\le C\varepsilon^2$.
The two $w$ boundaries point inward, because $z\le-q_g$ and
$h_s'=O(\varepsilon^3)$. The tube contains the entry and reaches
$z=-q_g$ in finite time. Its transition is strictly monotone for
each positive $\varepsilon$; no uniform lower derivative is required.

Choose a small rectangle $[-q_g,\xi]\times(0,\eta)$ with
$F_y=1+bz-2y>0$. Its top and the part of its right side outside
the outgoing section core have positive distance from the graphic.
At the left side $F>0$. A counted cycle cannot leave through the
bottom, left, or top, and must leave in finite time to return to
$\Sigma_u$. It therefore exits at the prescribed right section.
To exclude a zero of $F$ before that exit, use the exact identities
\[
 \dot F=F_zF+F_yG,\qquad \dot G=yF+zG.
\]
The quadrant $F\le0,G\le0$ is positively invariant in the
rectangle. At a first zero of $F$ reached from $F>0$, one has
$G\le0$; equality would give an equilibrium. Entering that quadrant
prevents the required right exit. Thus $F>0$ on the whole arc.

Finally, use $z$ as independent variable there. Between two passing
solutions of $y'=G/F$, scalar uniqueness orders every intermediate
solution. It stays between their graphs and has positive speed,
since $F_y>0$. Its target remains in the interval between their
targets. The domain of such passages is therefore an open interval.
Taking its inverse image under the monotone preceding transition
gives the asserted interval of $u_{\rm up}$. Equality
$T_+(u_{\rm up})=T_-(u_{\rm up})$ joins
two finite pieces of the same orbit; conversely every counted cycle
gives that equality. Uniqueness identifies distinct zeros with
distinct periodic orbits.
\end{proof}

For uniform regional estimates, fix the finite output order, the
physical sections and enlarged cores first. Then choose the selected
reference orbits, fibers and saddle collars, followed by compact
coefficient blocks, root gaps, barriers and tube widths. The final
parameter and Hausdorff shrinkings respect these choices; the actual
finite-passage interval is established before any return comparison.

\subsection{An explicit graphic-anchored annulus}
\label{sec:graphic-anchored-annulus}
Fix a compact interval $J\Subset(0,2)$ for this geometric construction.
Put $o=(0,1)$, $v=y-1$, and use the divided fast field
$z'=bz-v$, $v'=z$ only for this geometric construction.
Along its nonzero solutions the continuous polar angle about $o$
satisfies
\[
 \theta'=\frac{z^2-bzv+v^2}{z^2+v^2}
                         \ge1-\frac{\max J}{2}>0 .
\]
The explicit fast orbit starts with angle $-\pi/2$, passes
through the upper $z=0$ point, and reaches the foot with angle
$\theta_f(b)\in(\pi,3\pi/2)$. Its first and second parts have
$z>0$ and $z<0$, respectively, so it has not made an additional
turn. Along the axis segment from the foot to the origin the
angle continues strictly upward to $3\pi/2$, since
$\det((z,-1),(1,0))=1$.
Thus the whole graphic is a single positive radial graph
\[
 g_b(\theta)=o+R_b(\theta)e_\theta,\qquad
 e_\theta=(\cos\theta,\sin\theta).
\]
On the fast part its radius is obtained by the strictly monotone
angle-to-time inverse. On the axis part it is $-1/\sin\theta$.
The two expressions agree at the foot and at the circular
endpoint. The function $R_b(\theta)$ is jointly continuous,
piecewise smooth, and has common positive lower and upper
bounds on $J$. No derivative matching is claimed at the foot.

Choose $0<d<\min R_b/2$, also smaller than the upper-source core.
The explicit map
\[
 \Phi_b(\theta,t)=o+(R_b(\theta)+t)e_\theta,\qquad |t|\le d
\]
is globally injective. Different angles give different nonzero
radial rays and each radius is strictly increasing in $t$.
Its continuous inverse is
$\theta=\arg(x-o)$, $t=|x-o|-R_b(\arg(x-o))$.
It therefore embeds the entire cylinder, including both graphic
endpoints. Its fibers are anchored exactly on $\Gamma_b$ and
$|\Phi_b(\theta,t)-g_b(\theta)|=|t|$.
At $\theta=\pi/2$ the fiber is precisely $(0,y_b+t)$,
the original vertical source, whose transverse speed has a
fixed sign. Only joint continuity of $R_b$ is needed, so the radial
construction also applies at the singular endpoints.

Let $\mathcal A_b(r)=\Phi_b(S^1\times(-r,r))$.
An essential curve in this annulus meets every fiber: otherwise
cutting at a missed fiber puts it in a rectangle and makes it
contractible. Hence it is within two-sided Hausdorff distance
$r$ of the graphic. Conversely, for any fixed smaller annulus, joint
continuity of the inverse gives a common Hausdorff radius placing every
cycle within that distance of the graphic in the annulus.
Lemma~\ref{lem:sections} makes it
cross the upper fiber exactly once. A contractible Jordan curve
would leave both ends of this crosscut in its exterior and have
an even number of transverse intersections. Thus the cycle is
essential.

For two nested essential counted cycles in the smaller annulus, their
intervening closed annulus remains there. The inner boundary
connects to $o$ without meeting either cycle and lies inside both;
the outer boundary connects to infinity and lies outside both.
Choose $\mathcal A_-\Subset\mathcal A_+$ with the larger annulus
within the phase and Hausdorff bounds of the regional multiplier
estimate, and place the final counted cycles in the smaller one.
Their disks and intervening annuli are in $y>0$, since a Jordan
disk's lowest ordinate is attained on its boundary.

\subsection{An annulus argument}
When attraction can be proved over a whole period, we can count cycles
without constructing a return map across every intervening separatrix.
The annulus argument below turns that multiplier estimate into a bound in
terms of the number of equilibria.

\begin{lemma}\label{lem:attracting-annulus}
Suppose a planar $C^1$ field has at most $M$ equilibria in an annular
neighborhood, and every periodic orbit essential in that annulus is
hyperbolic attracting. Then any collection of such cycles whose
intervening closed annuli remain in the neighborhood has at most
$M+1$ members.
\end{lemma}
\begin{proof}
Take any finite collection of nested cycles and let $E$ be the closed
annulus between two consecutive cycles. Suppose $E$ has no equilibrium.
Both boundary cycles are invariant; uniqueness prevents crossing them,
and compact continuation gives a complete two-sided flow on $E$.
The basin of the inner cycle relative to $E$ is nonempty and open,
while an attracting neighborhood of the outer cycle is disjoint from it.
Its nonempty relative boundary is compact, invariant, and separated
from both cycles. The omega-limit set of a point on this boundary is
nonempty and contains no equilibrium. The Poincar\'e--Bendixson theorem
\cite[Section~7.3, Lemma~7.13]{Teschl2012} makes it a periodic orbit.
If that orbit were contractible in the original annular neighborhood,
its Jordan disk would lie inside $E$. Its index is one, whereas a
zero-free disk has index zero
\cite[Section~3.12, Theorems~1--3]{Perko2001}. Thus it is essential
and, by hypothesis, attracting. The chosen boundary point belongs
to its open basin, which is disjoint from the inner cycle's basin.
This contradicts the definition of that boundary. Each intervening
annulus therefore contains an equilibrium.
There are at most $M$ such disjoint annuli. Applying the argument to
any finite subcollection proves the result also against a proposed
infinite collection.
\end{proof}

\begin{lemma}\label{lem:generic-attracting-index}
Let $X_p$ be a planar family continuous in the $C^1$ topology at $p_0$.
Suppose the counted cycles of $X_{p_0}$ are hyperbolic attracting, lie
in $y>0$, and are essential in an annular neighborhood
$\mathcal A_-\Subset\mathcal A_+$, with their intervening closed annuli
contained in $\mathcal A_-$.  Suppose there are admissible parameters
$p_k\to p_0$ such that every periodic orbit of $X_{p_k}$ contained in
$\mathcal A_+\cap\{y>0\}$ and essential in $\mathcal A_+$ is hyperbolic
attracting, and $X_{p_k}$ has
at most four equilibria in $y>0$, all nondegenerate.
Then there are at most two counted cycles of $X_{p_0}$.
The admissible sequence may lie in a one-sided parameter set.
\end{lemma}
\begin{proof}
Fix a finite collection of $n$ counted cycles and disjoint compact tubes
in $\mathcal A_-\cap\{y>0\}$. Their finite minimal returns have
$P'-1\ne0$, so hyperbolic persistence gives one nearby cycle in each
tube \cite[Section~12.2, Lemmas~12.6--12.7]{Teschl2012}.
Transverse end hits and separated middle arcs preserve minimality.
The persisted cycles remain positive, distinct, attracting, and nested;
their intervening annuli remain in $\mathcal A_+$ and their bounded
interiors in $y>0$. The required perturbation neighborhood may depend
on the finite collection. For all sufficiently large $k$, apply the
following argument to $X_{p_k}$.
The vector-field index along a periodic orbit in the positive geometric
orientation is $+1$, whether the flow traverses it clockwise or
counterclockwise.  The innermost disk therefore contains an equilibrium.
Lemma~\ref{lem:attracting-annulus}, applied to each intervening annulus
using the attracting hypothesis for $X_{p_k}$, shows that it too contains
an equilibrium.  Its index sum is $1-1=0$.  A nonempty collection of
nondegenerate equilibria with zero index sum has at least two members,
because the index of each member is $+1$ or $-1$.  All these disjoint
regions lie in $y>0$, where the total equilibrium bound is four.  Hence
\[
 1+2(n-1)\le4,
\]
and $n\le2$.  The argument for arbitrary finite collections excludes an
infinite collection as well.
\end{proof}

For the present polynomial family the nondegenerate parameters required
in this lemma form a dense set.  When
$\delta=\mu_4\ne0$, equilibria correspond to roots of
\[
 Q(z)=z^2(\mu_1+\mu_2z+\mu_3z^2)-\delta z-b\delta z^2-\delta^2,
\]
and $Q'(z)=z\det D(F,G)$ at an equilibrium. The multiple-root condition
is a proper algebraic condition in the four raw parameters. One exact
witness, valid for each fixed $b$, is
\[
 \kappa=6/625,\quad\delta=12/25,\quad
 \mu_3=-6/625,\quad\mu_2=-12/125,\quad
 \mu_1=(12/25)b-42/125 .
\]
It gives $Q=-\kappa(z+1)(z+2)(z+3)(z+4)$ and discriminant
$144\kappa^6>0$. A nonzero real polynomial cannot vanish on an open
box, by the one-variable polynomial identity theorem applied
successively to its coefficients. Its nonzero set is therefore dense
in every open raw window. The witness proves properness and need not
belong to the small parameter window. If a multiplier estimate holds
only for $\delta\ge0$, at a field with $\delta=0$ choose perturbations
with $\delta>0$: every coefficient neighborhood has a nonempty open
part on that side, which the algebraic exceptional set cannot contain.
For each fixed positive $\varepsilon$, the map from
$(E_0,E_1,E_2,D)$ to the four raw coefficients is an invertible diagonal
linear map.  Thus the generic perturbations can be chosen within the same
coefficient chart before radial normalization.  The attracting hypothesis
on $\mathcal A_+$ must be supplied by the regional multiplier estimate;
it does not follow from persistence of the original finite collection.

\begin{lemma}\label{lem:equilibria}
For sufficiently small raw parameters, system~\eqref{eq:family}
has at most four equilibria in $y>0$.
\end{lemma}
\begin{proof}
Write $\delta=\mu_4$. If $\delta\ne0$, then $y=-\delta/z$
and $z\ne0$, so the abscissa satisfies
\[
 z^2(\mu_1+\mu_2z+\mu_3z^2)
                    -\delta z-b\delta z^2-\delta^2=0.
\]
This is a nonzero polynomial of degree at most four. Each root
determines at most one positive ordinate. If $\delta=0$, a positive
equilibrium has $z=0$ and $y-y^2+\mu_1=0$, giving at most two.
\end{proof}

By Lemma~\ref{lem:sections}, every counted cycle lies in $y>0$.
Its Jordan disk also lies there, since the lowest ordinate of a bounded
Jordan disk is attained on its boundary. Thus the intervening annuli lie
in $y>0$, and Lemma~\ref{lem:equilibria} supplies the equilibrium bound
used with Lemma~\ref{lem:attracting-annulus}.
